\documentclass[10pt,a4paper]{amsart}
\usepackage[margin=19mm]{geometry}
\usepackage{amsmath,amssymb,mathtools}
\usepackage[scr=rsfs,cal=euler]{mathalfa}
\usepackage{xfrac}
\usepackage[unicode,hidelinks,bookmarks=false]{hyperref}
\newcommand{\ie}{i.e.\ }
\newcommand{\cf}{cf.\ }
\newcommand{\resp}{respectively\ }
\newcommand{\Subsection}{\subsection}
\providecommand{\nobreakdash}{\nobreak\hbox{-}}
\setlength{\emergencystretch}{2em}
\multlinegap0pt
\usepackage{amssymb}
\usepackage{enumerate}
\usepackage{tikz}
\usepackage[shortlabels]{enumitem}
\newtheorem{theorem}{Theorem}
\newtheorem{proposition}{Proposition}
\newcommand{\Z}[0]{{\mathbb{Z}}}
\def\lukas{\color{blue}}
\title{Localized frames without inequalities}
\author{Peter Balazs, Lukas K\"ohldorfer, Michael Speckbacher}
\date{Source: arXiv:2506.02862v3 (CC BY 4.0)}
\begin{document}
\maketitle

\noindent\emph{Source and licence.} Localized frames without inequalities. Version arXiv:2506.02862v3 (CC BY 4.0), distributed by arXiv under the Creative Commons Attribution 4.0 International licence, http://creativecommons.org/licenses/by/4.0/, as recorded in the arXiv licence metadata for this exact version and shown on the abstract page https://arxiv.org/abs/2506.02862v3. Full licence notice: \texttt{licenses/arxiv-2506.02862-CC-BY-4.0.txt}.\par\smallskip

\noindent\textit{Editorial scope.} Counted unit: the article's proposition beta (source0.tex lines 1410--1411). The article's complete original proof of it is delivered with it (source0.tex lines 1413--1428, one \texttt{proof} environment of 16 lines). No mathematics is written, repaired or re-derived: the statement and the proof are the article's own lines, in the article's own notation, and no retained line is substituted or re-flowed.  Retained uncounted as the source-local context the proof needs: lines 1362--1364; lines 1372--1383; lines 1385--1394.  Nothing was omitted from any retained span, so the package carries no omission record.  No retained line contains a citation, so the wrapper carries no bibliography and no bibliography entry.  No retained line contains a figure environment, an image-inclusion command or drawing code, so no image is delivered and the package declares no asset dependency.  Editorial formatting only, in addition: an AMS \texttt{amsart} preamble that restates the article's own declarations and macros used by the retained lines, namely the environments \texttt{theorem}, \texttt{proposition} on separate counters, together with the standard packages those lines need.  The retained environments print in this document as Theorem 1, Proposition 1 in order of appearance, whereas the article prints the same items with the numbering of its own full document.  The complete native source of the article as distributed in the arXiv e-print for this exact version is bundled unchanged as \texttt{sources/179230/source0.tex}.
\begin{equation}\label{pert}
\vert x_k -k\vert \leq C,\qquad\text{for some }C>0\text{ and every }k\in\Z^d. 
\end{equation}

Before doing so, we compute the R-dual $\omega$ of $\psi$. Since the frame operator $S_\varphi$ commutes with translations by $k\in \Z^d$, so does $S_\varphi^{-1/2}$, and we see that  $S_\varphi^{-1/2} T_k\varphi=T_kS_\varphi^{-1/2} \varphi=T_k \gamma$ for $\gamma:=S_\varphi^{-1/2} \varphi$ and $k\in\Z^d$. Consequently,
\begin{equation}\label{eq:omega-SIS}
\omega_k=\sum_{l\in \Z^d}\langle \psi_l , T_k\varphi \rangle S_{\varphi}^{- {1}/{2}}T_l\varphi=\sum_{l\in \Z^d}  \overline{\varphi(x_l-k)}  T_l \gamma,\qquad k\in\Z^d.
\end{equation}
Note, that in case $X=(x_k)_{k\in\Z^d}= \Z^d$, the R-dual $\omega$ of $\psi$ is again a family of shifts of a single generator $\eta$ as
$$\omega_k=T_k\sum_{l\in \Z^d} \overline{ \varphi(l-k)}  T_{l-k} \gamma=T_k \sum_{l\in \Z^d}  \overline{\varphi(l)}  T_{l}\gamma=T_k\eta,\qquad k\in\Z^d.$$
Computing the entries of the associated Gram matrix $G_\omega$, we see that 
\begin{flalign}\label{correlationmatrix}
[G_\omega]_{k,n} &= \langle \omega_n, \omega_k \rangle = \left\langle  \sum_{l\in \Z^d}  \overline{\varphi(x_l-n)} S_{\varphi}^{-1/2} T_l \varphi , \sum_{m\in \Z^d}  \overline{\varphi(x_m-k)} S_{\varphi}^{-1/2} T_m \varphi \right\rangle \notag \\
&= \sum_{l\in \Z^d} \varphi(x_l-k) \overline{\varphi(x_l-n)},
\end{flalign}
since $(S_{\varphi}^{-1/2} T_l \varphi)_{l\in \Z^d}$ is an orthonormal basis. %{\lukas{Is this correlation matrix somewhat known?}}

\begin{theorem}\label{samplingthm}
Assume that $\varphi \in L^2(\mathbb{R}^d)$ satisfies the conditions (i)$-$(iii), and let $(x_k)_{k\in \mathbb{Z}^d} \subset \mathbb{R}^d$ be relatively separated. Let $G_\omega$ be the autocorrelation matrix with entries given by (\ref{correlationmatrix}) and assume that condition (\ref{pert}) holds. Then the following are equivalent.
\begin{itemize}
\item[(a)] $(x_k)_{k\in \mathbb{Z}^d}$ is a stable set of sampling for $V^2(\varphi)$.
\item[(b)] $(x_k)_{k\in \mathbb{Z}^d}$ is a stable set of sampling for $V^{\infty}(\varphi)$ and $D_{\psi}:\ell^{\infty} \to V^{\infty}(\varphi)$ has closed range. 
\item[(c)] $G_\omega$ is invertible on $\ell^1(\Z^d)$.
\item[(d)] $G_\omega$ is invertible on $\ell^{\infty}(\Z^d)$.
\item[(e)] $G_\omega$ is invertible on $\ell^{2}(\Z^d)$.
\end{itemize} 
\end{theorem}

\begin{proposition}\label{beta} $\mathbb{Z} - \tau$ is a stable set of sampling for $V^2(\beta)$ if and only if $\tau \neq \frac{1}{2}$.
\end{proposition}

\begin{proof} We check condition (e) of Theorem \ref{samplingthm}. Via (\ref{correlationmatrix}), we see that
$$[G_\omega]_{k,n} = \sum_{l\in \mathbb{Z}} (1-\vert l - \tau -k\vert)_+ \cdot (1-\vert l - \tau -n\vert)_+ = \begin{cases}
(1-\tau)^2 + \tau^2 \quad &(k=n)\\
(1-\tau)\tau \quad &(\vert k-n\vert = 1)\\
0 \qquad &\text{(else)}
\end{cases}$$
due to support reasons. Thus, $G_\omega$ is an infinite Toeplitz matrix, i.e., a convolution operator with respect to the sequence $g = (g_k)_{k\in \mathbb{Z}}$ with entries corresponding to the constant diagonals of $G_\omega$. %$(g_k)_{k\in \mathbb{Z}}$, where $g_0 = (1-\tau)^2 + \tau^2$, $g_1 = g_{-1} = (1-\tau)\tau$ and $g_k = 0$ for $\vert k \vert \geq 2$. 
By unitary equivalence via the Fourier transform on $\mathbb{Z}$, $G_\omega$ is invertible on $\ell^2(\mathbb{Z})$ if and only if $\widehat{g}$ is nowhere vanishing. We compute 
\begin{flalign}
\widehat{g}(\xi) &= (1-\tau)\tau e^{2\pi i \xi} + (1-\tau)^2 + \tau^2 + (1-\tau)\tau e^{- 2\pi i \xi} \notag \\
&= (1-\tau)^2 + \tau^2 + 2(1-\tau)\tau \cos(2\pi \xi) \notag \\
&\geq (1-\tau)^2 + \tau^2 - 2(1-\tau)\tau \notag \\
&= (2\tau - 1)^2, \notag
\end{flalign}
and see that $\widehat{g}$ is non-vanishing precisely when $\tau \neq \frac{1}{2}$.
\end{proof} 

\end{document}
